\chapter{Outlook: Cubilox Mk2}
\label{ch:outlook}

The experience from the Mk1 leads to the following plans for a Mk2:

\begin{itemize}
  \item \textbf{A real printed circuit board} with proper connectors instead of the prototype board
        and soldered jumper wires. This saves space and weight and makes the wiring more reliable.
  \item \textbf{Stronger motors.} More torque would allow calmer balancing with more reserve and
        could make it possible for the cube to jump up from a resting position, like the
        Cubli~\cite{gajamohan2012cubli}.
  \item \textbf{An RGB status LED.} Start-up, calibration and state changes would be shown by light;
        the buzzer would only be used for important events such as errors or an empty battery.
\end{itemize}

Further possible improvements that follow from the limitations in \cref{ch:limitations}:
\begin{itemize}
  \item \textbf{Safe battery connection} with insulated contacts.
  \item \textbf{Model-based control.} With an identified model of the cube, the gains could be
        computed (for example with a linear-quadratic regulator) instead of being tuned only by
        experiment.
\end{itemize}
